\section{AI-agent usage in this analysis program}
\label{sec:ai_agent}

This work includes an explicit evaluation of AI-agent support for an
experimental QCD analysis workflow. The central question is not only whether an
agent can generate text or isolated scripts, but whether it can execute a full
analysis loop under physicist supervision: data and Monte Carlo (MC) processing, unfolding
orchestration, systematic propagation, figure production, and note maintenance
with traceable links to code and outputs.

\subsection{Scope and operating model}

The agents used in this analysis are Anthropic Claude Code and OpenAI Codex, depending on the preference of the physicist.
The operating model is human-in-the-loop: the physicist formulates goals and
provides expert direction through iterative prompts; the agent executes the
corresponding implementation steps; the physicist reviews the outputs and
provides feedback; the process is then repeated. Physics decisions remain
analyst-owned, while the agent provides implementation speed, consistency
checks, and documentation rigor. In practical terms, the agent is used to:
\begin{itemize}
\item write, execute, and monitor analysis code,
\item enforce consistent plotting and note formatting conventions for work communication,
\item maintenance an understanding of which analysis code and configurations,
\item iterate quickly on ideas other time consuming tasks.
\end{itemize}

\noindent Our operating model is intentionally conservative to see how the agent performs when used to accelerate the common physics engineering and documentation tasks, rather than the core scientific decision tasks.

\subsection{Inputs and context supplied to agents}
\label{sec:agent_inputs}

The context supplied to the agent at the start of the session are the following.
\begin{itemize}
\item Context: the published ALEPH thrust measurement, prior ROOT-based analysis macros for event selection and thrust calculation, and internal analysis notes documenting the open data format, detector calibration, and object definitions~\cite{Heister:2003aj,ALEPH:docs}.
\item Datasets: the ALEPH LEP1 data files for the 1994 running period at $\sqrt{s}=91.2$~GeV, together with the corresponding ALEPH MC samples. Additional modern MC were given for theoretical uncertainty estimation.
\item Physicist-written prompts: prompts specifying the physics objective (e.g.\ ``propagate the hadronic-event correction as an independent systematic component with per-bin envelope uncertainties''), acceptance criteria (e.g.\ ``refolded spectra must reproduce detector-level data within statistical precision''), and output-format expectations (figure style, table layout, note section to update).
\end{itemize}

\subsection{Boundaries, controls, and reproducibility}

To keep scientific accountability explicit, each high-impact step has a control point where analyst sign-off is required before results are treated as baseline. The analysis stages where control points were used are: event and object selection definition, unfolding regularization choice, systematic-combination policy, final figure and uncertainty-presentation choices, interpretation description of the work. Beyond these sign-off gates, the physicist performed concrete validation checks at each major analysis stage by inspecting the artifacts of the agents work. These checks are enumerated here.

\begin{itemize}
    \item Event selection: Cut-flow yields were compared to the published ALEPH numbers~\cite{Heister:2003aj} to verify selection efficiency. Distributions of $\cos\theta_{\mathrm{thrust}}$ and $\varphi$ were inspected for detector artifacts; this step uncovered a localized $\varphi$ excess in the data that was traced to a detector region and checked before proceeding.
    \item Unfolding: Refolded spectra were required to reproduce the detector-level data within statistical precision. An iteration scan confirmed that the refolded-to-measured is stable as the regularization parameter is scanned.
    \item Systematic uncertainties: Each variation's impact on the unfolded distribution was inspected for large bin-to-bin fluctuations. The covariance and correlation matrices for individual components as well as the total were computed.
    \item Results: The unfolded thrust distribution was compared to the ALEPH\,2004 published result over the 42~reported bins. Per-bin uncertainty breakdowns and correlation matrices are reported.
    \item All figure captions, table entries, and interpretive statements in this work were reviewed by the physicist for physics accuracy. Agent-generated text was treated as a draft requiring human approval before inclusion in the baseline.
\end{itemize}

Reproducibility is ensured through multiple complementary mechanisms. All analysis steps are version-controlled, with code and configuration managed in a git repository so that any result can be traced to a specific commit and set of inputs. Workflow runs produce structured output artifacts, including summary files that capture the configuration used, the software versions, and caches of the exported spectra and covariance files. Physicist oversight at each review stage provides a layer of verification, with each step requiring explicit physicist approval before being promoted to the nominal result and the analysis to continue. These controls ensure that results are reproducible and that the physicist retains accountability for the scientific claims.
